\chapter{Insider threats}

DPRK IT workers are fraudulently hired remote staff who
fund a sanctioned regime. They can also steal, extort,
or sabotage the supply chain. Hiring and access control
are the primary defenses, not antivirus after the fact.

This is not a rare edge case. It is a hiring-process
failure mode with legal and theft consequences.

\section{Insights}

\begin{itemize}
  \item Fake identity is the job. Interviews, laptops,
  and stand-ins are part of the method.
  \item Shared credentials and laptop-from-home setups
  are how one hire becomes the whole repo.
  \item Paying the wrong person can be a sanctions
  problem, not only an HR problem.
  \item After discovery, assume they mapped internals.
  Rotate, do not just fire.
\end{itemize}

\section{Quick wins}

\begin{enumerate}
  \item Video interviews with identity that matches
  payroll. Refuse stand-ins.
  \item Company-owned machines. No personal laptops on
  production.
  \item Least privilege from day one. No standing admin
  for contractors.
  \item If a hire already feels off, read the DPRK
  pages and the incident playbook before you ping them
  in Slack.
\end{enumerate}
